\section{INTRODUCCIÓN}

\subsection{Importancia del Modelado y Simulación de \ldots}

\subsection{Trabajos previos en Modelado y Simulación Dinámica de \ldots}

\subsection{Planteamiento del Problema a resolver}

\subsection{Objetivos}

\begin{table}[htbp]
	\caption{Propiedades mecánicas del aluminio 6061-T6} % Solo escribe el título normal
	\label{tab:propiedades}
	\centering
	\begin{tabular}{lcc}
		\toprule % Línea superior
		Propiedad & Valor & Unidad \\
		\midrule % Línea debajo de los encabezados
		Módulo de elasticidad & 68.9 & GPa \\
		Límite de fluencia & 276 & MPa \\
		Resistencia última & 310 & MPa \\
		Densidad & 2700 & kg/m$^3$ \\
		\bottomrule % Línea inferior final
	\end{tabular}
	\vspace{0.2cm}
	
	\raggedright % La nota va alineada a la izquierda
	\textit{Nota.} Datos obtenidos de Callister y Rethwisch (2020).
\end{table}

\begin{figure}[htbp]
	\caption{Modelo CAD de la cortadora láser CNC}
	\label{fig:cad}
	\centering
	\includegraphics[width=0.8\textwidth]{example-image-a}
	\vspace{0.2cm}
	
	\raggedright
	\textit{Nota.} Diseño preliminar del equipo ensamblado.
\end{figure}